
\appendix
\section{Appendix}\label{sec:appendix}

\subsection{Proof of Lemma~\ref{lem:finiteness} -- Finiteness of nodes}


\begin{proof}
Let $n=|\varphi|$ and let $M$ be the largest constant in any timing interval of $\varphi$. Let $Sub(\varphi)$ be the Fischer--Ladner closure of $\varphi$: the least set containing $\varphi$ and closed under subformulas and under the one-step unfoldings the rules produce -- the released form $\psi\land(\varphi\lor\AX r)$, the marker-decorated formulas $\EX e$, $\AX e$, and the pursued forms $\xi^{\bullet}$, $\EX\xi^{\bullet}$, $\AX\xi^{\bullet}$ of timed untils. Each unfolding only combines existing subformulas with a single marker, so the closure is linear: $|Sub(\varphi)|=O(n)$.

By construction every formula in $\Gamma$ or $\mathcal{X}$ belongs to $Sub(\varphi)$, so $\Gamma$ and $\mathcal{X}$ range over $2^{O(n)}$ subsets each; likewise $\mathcal{E}$ ranges over the eventualities of $Sub(\varphi)$, $2^{O(n)}$ values.

For clocks we use one {formula-clock} per timed-until subformula (Section~\ref{sec:tableau}). A timed until demanded by a persistent context is verified {independently} from each state, on its own dedicated-check branch~\eqref{tag:chk}, with its clock reset there; the dense family of per-state demands is never accumulated as concurrent obligations, and along any single branch at most the timed nesting depth of clocks is live. Over the whole tableau there are therefore $K=O(n)$ clocks, and the map $\mathcal{O}$ contributes $2^{O(n\log n)}$ values.

Over $K=O(n)$ clocks bounded by $M$, the number of canonical extrapolated DBMs is the number of $(K{+}1)^2$ matrices whose entries are bounds in $(\{-M,\dots,M\}\times\{<,\le\})\cup\{(\infty,<)\}$, i.e.\ $M^{O(n^2)}$ \cite{Bouyer05TimedAutomata,AlurDill1994}. This is the count of canonical {zones} (DBMs), correcting the region count $O(K!\,(2M{+}2)^K)$, which over-counts by tracking orderings of fractional parts irrelevant to a canonical DBM.

Taking the product, the tuple $\langle\Gamma,\mathcal{X},Z,\mathcal{O},\mathcal{E}\rangle$ ranges over at most $2^{O(n)}\cdot 2^{O(n\log n)}\cdot M^{O(n^2)}=2^{O(n\log n)}\cdot M^{O(n^2)}$ values. This is finite and, with timing constants in unary, singly exponential in $|\varphi|$. The dense set of per-state until-demands under persistent contexts collapses to these finitely many $(\text{node},\text{region})$ classes -- the region quotient at the heart of the decidability of the non-punctual fragment~\cite{fsttcs2000-LN,LaTorreNapoli2001TimedTreeAutomata}.


%%Let $\Sub\varphi$ be the closure of $\varphi$. 
%%By construction of the rules, every formula that can appear in
%%$\Gamma$ or $\mathcal{X}$ belongs to $\Sub\varphi$.
%%Hence $\Gamma$ and $\mathcal{X}$ range over finitely many subsets of
%%$\Sub\varphi$. Similarly, $\eventualities$ ranges over subsets of the finitely many
%%eventualities in $\Sub\varphi$
%%
%%
%%Clock allocation introduces at most one clock per timed eventuality
%%in $\Sub\varphi$, therefore the set of clocks is finite. 
%%$\obligations$ is a map from timed eventualities to clocks, i.e. between two finite sets, hence finite. Over a fixed finite clock set, the abstraction $\alpha_M$ yields only
%%finitely many abstract zones. Since zones are stored in canonical form
%%modulo $\cong_M$, there are only finitely many possible values of $Z$. Therefore the tuple $\langle \Gamma,\mathcal{X},Z,\obligations,\eventualities\rangle$ can take only finitely many distinct values, and only finitely many distinct nodes can arise.
%%
\end{proof}


\subsection{Proof of Theorem~\ref{thm:termination} -- Termination of the tableau procedure}


\begin{proof}
Fix an input formula $\varphi$ and consider any branch $N_0,N_1,\dots$ generated by the depth-first construction.
By Lemma~\ref{lem:finiteness}, only finitely many distinct tableau nodes can arise during the construction. In particular, an infinite branch would necessarily contain a repeated node by the pigeonhole principle.

If the branch contains no repeated node, then its length is bounded by the number of distinct tableau nodes and is therefore finite. Otherwise, there exist indices $i<j$ such that $N_i \equiv N_j$.
When such a repetition is encountered, the loop rule is applied,
and the branch is either declared successful or closed.
In either case, no further unfolding beyond $N_j$ takes place.

Thus every branch is finite. Each node has finite out-degree -- at most five children from an expansion rule and at most $|\Gamma|=O(n)$ from successor generation -- so by König's lemma a construction all of whose branches are finite is itself finite; equivalently, the memoised procedure expands each of the finitely many distinct nodes (Lemma~\ref{lem:finiteness}) exactly once. Hence the tableau procedure terminates.
\end{proof}


\subsection{Proof of Lemma~\ref{lem:nonzeno} -- Non-Zeno characterisation}
\begin{proof}
($\Rightarrow$, necessity of time progress) Suppose the loop is {not} time-progressing. Either (a) it can elapse only zero time, so every unfolding has $\sum_k t_k=0$ and is Zeno; or (b) some clock $x$ is never reset on the segment yet, on it, keeps a finite upper bound $c\le M$ (it is not extrapolated above $M$). Since $x$ is never reset, its value is non-decreasing along the $\omega$-unfolding and after $m$ traversals equals $x_{\text{entry}}+\sum_{k} t_k$ accumulated so far; the constraint $x\le c$ then forces $\sum_k t_k\le c-x_{\text{entry}}<\infty$, so the unfolding is Zeno. Hence a non-Zeno unfolding requires time progress. (This is exactly the $\mathrm{EG}\,p\land\mathrm{AG}_{[1,M]}\neg p$ loop, whose release clock $x_\rho<1$ is never reset.)
($\Leftarrow$, sufficiency) Suppose the loop is time-progressing and fix a delay profile realising total delay $\delta>0$ on one traversal. Because $N_i\equiv N_j$, the segment repeats structurally each iteration; by hypothesis every clock that bounds the delay from above is either reset within the segment -- so its value at the start of each iteration equals its value on the first, replenishing its budget -- or extrapolated above $M$, hence imposing no finite upper bound. Therefore the same profile is realisable on every iteration, and the unfolding accumulates $\sum_k t_k=\sum_{\mathrm{iter}}\delta=\infty$: it is non-Zeno, with every iteration elapsing the same positive amount.
\end{proof}


\subsection{Zone-path feasibility}
\begin{lemma}[Zone-path feasibility]\label{lem:feasible}
Let $N_0,N_1,\dots$ be an infinite path of a marking. Write $Z_i^{\flat}$ for the {exact} reachable zone along the path: $Z_0^{\flat}$ is the initial zone and
\[
Z_{i+1}^{\flat}=\big(((\uparrow Z_i^{\flat})\cap\mathsf{Inv}(s_i))\cap g_i\big)[R_i:=0]
\]
is the true post-image under the delay $\uparrow$ (restricted to the location invariant $\mathsf{Inv}(s_i)$), the guard $g_i$ (a single-clock bound $x\in\kI$ or $\top$), and the reset $R_i$. The procedure instead stores the {iterated extrapolation} $W_0=\mathrm{canonical}(\alpha_M(Z_0^{\flat}))$, $W_{i+1}=\mathrm{canonical}(\alpha_M(\mathrm{post}_i(W_i)))$, used only for loop detection; these over-approximate ($Z_i^{\flat}\subseteq W_i$), but because all guards are diagonal-free, iterated max-constant forward analysis is sound and {step-wise complete}~\cite{Bouyer05TimedAutomata} -- the per-edge-sequence biconditional -- so $W_i\neq\emptyset \iff Z_i^{\flat}\neq\emptyset$ along the path. The location invariant $\mathsf{Inv}(s_i)$ written in the recurrence is not an extra constraint the procedure lacks: on an open marking the stored zone already entails it -- the release entry caps $x_\rho<l_\rho$ (resp.\ $\le l_\rho$ for a left-open window) are exactly the T-ER before-region constraint, and the until deadline caps $x_\xi\le u_\xi$ are enforced by the integer-hyperplane split and deadline-closure -- and since all invariant constants are $\le M$, $W_i\models\mathsf{Inv}(s_i)\iff Z_i^{\flat}\models\mathsf{Inv}(s_i)$; so $\cap\,\mathsf{Inv}(s_i)$ is consistent with, not additional to, the invariant-free zone the loop test compares. Then there exist valuations $\nu_i\in Z_i^{\flat}$ and delays $t_i\ge 0$ with $\nu_i+t_i\in\mathsf{Inv}(s_i)\cap g_i$ and $\nu_{i+1}=(\nu_i+t_i)[R_i:=0]$, forming a concrete run respecting the invariants; on an accepting lasso the $t_i$ have strictly positive total delay per loop iteration.
\end{lemma}
\begin{proof}
Since the marking is open, the stored $W_i$ are non-empty, hence by forward-analysis completeness the exact $Z_i^{\flat}$ are non-empty. By the post-image identity every $\nu_{i+1}\in Z_{i+1}^{\flat}$ has a predecessor $\nu_i\in Z_i^{\flat}$ and a delay $t_i$ with $\nu_i+t_i\in\mathsf{Inv}(s_i)\cap g_i$ and $\nu_{i+1}=(\nu_i+t_i)[R_i:=0]$; picking any $\nu_{i+1}\in Z_{i+1}^{\flat}$ and threading backwards realises a finite prefix from the {exact} zones -- the valuations added by extrapolation play no role. For an accepting lasso the loop closes at $N_j\equiv N_i$, i.e.\ $W_j=W_i$; we do {not} infer any equality of the exact zones, and build the periodic run from the {reset structure} instead. By the time-progress condition (Lemma~\ref{lem:nonzeno}) every clock the loop bounds below $M$ is reset within $\pi_{\mathsf{loop}}$, so after one traversal those clocks return to identical concrete values; the remaining clocks lie above $M$, where no single-clock guard distinguishes two values, so their guard-outcomes are invariant. (An open loop cannot guard a non-reset clock by an upper bound $\le M$ -- that branch would close -- so only lower bounds $x\ge k$, satisfied as the clock grows, touch them; the run stays valid through warm-up and steady state.) Hence the concrete configuration recurs {modulo} the values of the clocks above $M$, and elapsing the same $\delta>0$ each iteration yields a non-Zeno run that is periodic in its guard-relevant projection (the concrete clocks above $M$ drift upward; a finite warm-up may let non-reset clocks pass $M$). Threading the prefix backward from the loop entry gives $\nu_0\in Z_0^{\flat}$. This is the standard concretisation of a diagonal-free zone cycle into an infinite timed run.
\end{proof}


\subsection{Proof of Lemma~\ref{lem:lasso} -- Lasso extraction from successful markings}
\begin{proof}
Let $\mathcal{M}$ be a successful marking. Since the tableau has finitely many distinct nodes (Lemma~\ref{lem:finiteness}) and $\mathcal{M}$ retains all AND-children and one OR-child per node, $\mathcal{M}$ is a finite graph. Consider any infinite path $\pi$ of $\mathcal{M}$. By finiteness $\pi$ repeats a node, so $\pi=\pi_{\mathsf{pref}}\cdot\pi_{\mathsf{loop}}^\omega$ for the first repetition $N_i\equiv N_j$; the cycle is closed by the loop rule's back-edge, or, at a node with no existential demand, by the stutter successor that re-enters the same node (so there are no childless leaves). Because $\mathcal{M}$ is successful, every postponed eventuality is fulfilled on $\pi_{\mathsf{loop}}$ and $\pi_{\mathsf{loop}}$ is time-progressing, i.e.\ the cycle is accepting. As $\pi$ was arbitrary, every infinite path of $\mathcal{M}$ is an accepting lasso.
\end{proof}


\subsection{Proof of Theorem~\ref{thm:soundness} -- Soundness of the tableau procedure}

%Because the branch is open, the zone $Z$ is non-empty at every node. 
\begin{proof}
Let $\mathcal{M}$ be a successful marking (Theorem hypothesis, via Lemma~\ref{lem:lasso}). We construct a TTS $\mathcal{T}=\langle S,s_0,\mathsf{Inv},E,\mathcal{P},L\rangle$ (Definition~\ref{def:tts}) and show $\mathcal{T},(s_0,\nu_0)\models\varphi$.

{Locations and edges.} The locations $S$ are the saturated nodes of $\mathcal{M}$, $s_0$ the root, $\mathcal{P}$ the propositions of $Sub(\varphi)$, and $L(s)$ the positive literals of $\Gamma_s$ (consistent because $s$ is open). The discrete edges are the $\EX$-successors, but read off the {pursuit} branches, not the bookkeeping continuations: for each $\EX\psi_m\in\Gamma_s$ -- whether a top-level demand or one generated inside a dedicated check -- we add an edge $(s,R_m,s'_m)$ to the corresponding successor, which carries $\psi_m$, the whole context $\mathcal{X}_s$, the active formula-clocks, and, for a universal timed modality, the deadline/entry invariant. The marking retains {all} such successors, so a node demanding several existential successors, or imposing $\AX$-obligations, contributes all of these edges (the step the previous single-branch construction omitted). By contrast, the left child of a dedicated check~(\ref{tag:chk},~\ref{tag:chkr}), which {drops} the demanded modality, contributes {no run of its own}: it is bookkeeping that lets a persistent context re-demand the modality afresh at the next state (spawning a new check). Hence every run of $\mathcal{T}$ passes only through pursuit-committed locations, where $\psi/q$ is asserted and the urgency invariant holds; the witness a check exhibits is a genuine run of $\mathcal{T}$ from the demanding configuration. The left-drop nodes are {unreachable} in $\mathcal{T}$ -- no edge of $E$ targets them -- so they organise the search without contributing any run; the computed $\mathrm{Sat}(\cdot)$ then coincides with the semantic satisfaction set in $\mathcal{T}$ (Lemma~\ref{lem:sat}).

{Universal timed modalities, configuration by configuration.} A single formula-clock cannot, along one spine, host the continuum of overlapping live windows that a persistent universal-timed context demands; we therefore do {not} establish a universal timed $\xi$ by one live clock. Instead its check certifies $Z\subseteq\mathrm{Sat}(\xi)$ at the demanding node, with the deadline/entry invariant making $\mathrm{Sat}(\xi)$ the genuinely {universal} set (every run from a valuation in it meets the window, not merely one). Thus every reachable configuration $(s,\nu)$ of $\mathcal{T}$ has $\nu\in\mathrm{Sat}(\xi)$, so $\mathcal{T},(s,\nu)\models\xi$; and $\mathcal{T}\models\mathrm{AG}\,\xi$ follows from this per-configuration membership rather than from any spine invariant. This is the region-quotient enforcement of~\cite{fsttcs2000-LN,LaTorreNapoli2001TimedTreeAutomata}, internalised.

{Invariants.} We equip each location with the invariant that forces the urgency demanded by its active {universal} timed obligations, so that no non-Zeno run may slip across a region boundary by lingering. For a node $s$, $\mathsf{Inv}(s)$ is the conjunction of:
\begin{itemize}
\item $x_\xi \le u_\xi$ (resp.\ $<u_\xi$ if $\kI_\xi$ is right-open), for every universal until $\xi=\tAU{\varphi}{\psi}{\kI_\xi}$ active and unfulfilled at $s$ with finite upper bound $u_\xi$ -- its {deadline}; and
\item $x_\rho < l_\rho$ (resp.\ $\le l_\rho$ if $\kI_\rho$ is left-open), for every universal timed release $\rho=\tAR{\varphi}{\psi}{\kI_\rho}$ active at $s$ and still in its {before}-region $Z_s^{<\kI_\rho}$ with $l_\rho>0$ -- its {window-entry} bound;
\end{itemize}
and $\mathsf{Inv}(s)=\top$ if there is none. The until caps force the witness to be reached by the deadline, with strictness {matching} the interval (the witness may be taken at a closed $u_\xi$). The release entry caps force a run to leave the before-region -- into the in-region, where $\psi$ {is} asserted -- by the time its window opens; their strictness is the {complement} of the window's left endpoint, opposite to the until rule. The reason is the super-dense reading of Definition~\ref{def:sat}: for a left-{closed} window $[l_\rho,\cdot]$ the run must already be in the $\psi$-location {at} time $l_\rho$ (the pre-jump position at $l_\rho$ would otherwise be a $\neg\psi$ position inside the window), so the cap must be strict, $x_\rho<l_\rho$, forcing the step strictly before $l_\rho$; for a left-{open} window $(l_\rho,\cdot]$ the instant $l_\rho$ is outside the window, so the non-strict cap $x_\rho\le l_\rho$ is correct and a strict one would block the only entry. Symmetrically the run stays in the $\psi$-location until after the window's right end, so $\psi$ holds at every position with time in $\kI_\rho$; asserting $\psi$ on the entry neighbourhood before $l_\rho$ is harmless (the before-region permits, but does not require, $\psi$). This is the lower-bound (entry) dual of the until's upper-bound (deadline) urgency. Together with the boundary-split of Section~\ref{sec:tableau}, every universal timed obligation is thereby enforced region-by-region on {all} non-Zeno runs, not merely along the marking's discrete skeleton. Existential obligations impose no invariant (one witnessing run suffices). Non-blockingness is recorded separately.

\begin{lemma}[Non-blocking]\label{lem:nonblock}
The constructed $\mathcal{T}$ is non-blocking: from every configuration on the marking a non-Zeno run exists.
\end{lemma}
\begin{proof}
$\mathsf{Inv}(s)$ is satisfiable -- all clocks at $0$ meet every $\le/<$ cap -- and along the open marking each active universal-timed clock stays within its current cap (an until past its deadline, or a release whose entry was skipped, would make the node closeable), so $\nu_s\models\mathsf{Inv}(s)$. Successor generation always supplies an outgoing edge (a demanded $\EX$-child, or the stutter edge), whose target $s'$ lies on the open marking and hence has $\nu_{s'}\models\mathsf{Inv}(s')$; thus a discrete step is enabled at or before each invariant boundary. Interleaving these steps with the positive per-iteration delays of Lemma~\ref{lem:nonzeno} produces a time-divergent run; the caps stack consistently because along the marking all active universal clocks simultaneously respect their bounds.
\end{proof}

{Valuations and time.} By Lemma~\ref{lem:feasible} the non-empty exact zones along $\mathcal{M}$ admit a coherent assignment of valuations $\nu_s\in Z_s^{\flat}$ (each satisfying $\mathsf{Inv}(s)$, since the active universal-until clocks are kept within their deadlines along the marking) and delays realising each delay-then-reset edge; by condition~(ii) and Lemma~\ref{lem:nonzeno} every loop of $\mathcal{M}$ is time-progressing, so the runs we extract along $\mathcal{M}$ are time-divergent and $\mathcal{T}$ is non-blocking. ($\mathcal{T}$ also admits Zeno runs, e.g.\ all-zero delays; these are simply excluded by the semantics, which quantifies over $\mathrm{Runs}^\infty$.) The universal obligations below are nonetheless verified against {all} non-Zeno runs, which the location invariants constrain.


We prove, by induction on the structure of $\eta\in Sub(\varphi)$, the claim: for every node $s$ of $\mathcal{M}$, if $\eta\in\Gamma_s$ then $\mathcal{T},(s,\nu_s)\models\eta$. The base and Boolean cases are immediate: literals hold by definition of $L$, and $\land,\lor$ by saturation and the inductive hypothesis (the chosen disjunct lies in $\Gamma_s$).

{Markers.} If $\EX\psi\in\Gamma_s$, the construction adds an edge to a successor $s'$ with $\psi\in\Gamma_{s'}$; by the inductive hypothesis $(s',\nu_{s'})\models\psi$, witnessing the existential successor. If $\psi\in\mathcal{X}_s$ (the source of every $\AX\psi$ obligation), then $\psi$ is inserted into the obligation set of {every} successor child; by the inductive hypothesis every discrete successor of $s$ satisfies $\psi$, so the universal-next requirement holds across the branching -- not vacuously through a single successor.

{Until (least fixpoint), $e=\tEU{\varphi}{\psi}{\kI}$ and the $\mathbf{A}$-variant.} When $e$ is demanded at $s$ it is verified by its dedicated check~\eqref{tag:chk}, which succeeds only if $Z_s\subseteq\mathrm{Sat}(e)$; thus the zone retained at $s$ contains {only} valuations from which $e$ holds. Fix any $\nu\in Z_s$. By $\nu\in\mathrm{Sat}(e)$, the backward fixpoint exhibits a finite path $\pi$ of $\mathcal{M}$ from $s$ -- following the go branches of $e^{\bullet}$, each a discrete successor of $\mathcal{T}$ retained by $\mathcal{M}$ -- to a stop node where $\psi$ holds and the elapsed time, tracked by $x_e$, lies in $\kI$, with $\varphi$ at every earlier position and every ambient safety obligation of $\mathcal{X}$ respected (these were carried in the check). The check edges are edges of $\mathcal{T}$, so $\pi$ is a run of $\mathcal{T}$ from $(s,\nu)$; hence $\mathcal{T},(s,\nu)\models e$ for the existential variant, as required by $\exists\rho$ in Definition~\ref{def:sat}. For the universal $\mathbf{A}$-variant the fixpoint quantifies over {all} successors and the demanding location carries the deadline invariant $x_e\le u_e$ (resp.\ $<u_e$); so no run may delay past $u_e$, every run is forced through a stop node within $\kI$, and $\mathrm{Sat}(e)$ is then the set of valuations from which {every} run does so. Since $Z_s\subseteq\mathrm{Sat}(e)$, every $\nu\in Z_s$ satisfies $\forall\rho$ in Definition~\ref{def:sat}. (When $u_e=\infty$ no window can be overshot and the eventuality condition~(i) alone suffices.)

{Release (greatest fixpoint), $\rho=\tER{\varphi}{\psi}{\kI}$ and the $\mathbf{A}$-variant.} When $\rho$ is demanded at $s$ it is verified by its check~\eqref{tag:chkr}, which succeeds only if $Z_s\subseteq\mathrm{Sat}(\rho)$; the retained zone thus contains only valuations from which $\rho$ holds. $\mathrm{Sat}(\rho)$ is the greatest-fixpoint of the partition~\eqref{tag:ter} (read co-inductively, $x_\rho$ reset at $s$), and its in-region branch {asserts} $\psi$, committing $\psi$ to the model locations occupied during $\kI$. Fix $\nu\in Z_s\subseteq\mathrm{Sat}(\rho)$. Before the interval ($x_\rho\in\kI^{-}$) $\psi$ is not required and the obligation propagates or is released early by $\varphi$; inside ($x_\rho\in\kI$) the committed location asserts $\psi$ unless $\varphi$ has released; after ($x_\rho\in\kI^{+}$) the obligation has expired. For the {existential} $\rho$ this holds along the witnessing run; for the {universal} $\mathbf{A}$-variant it must hold on {every} non-Zeno run, which two mechanisms together secure. The boundary-split materialises the in-region location (asserting $\psi$), and the entry invariant ($x_\rho<l_\rho$ for a left-closed window, $x_\rho\le l_\rho$ for a left-open one -- the complement of the window's left endpoint) forbids any run from delaying into the window at a location not yet asserting $\psi$: for a left-closed window the strict cap forces the step {strictly before} $l_\rho$, so the run is in the $\psi$-location {at} time $l_\rho$ and the pre-jump position there is not a $\neg\psi$ position inside the window; for a left-open window $l_\rho$ is outside the window and the non-strict cap suffices. As $Z_s\subseteq\mathrm{Sat}(\rho)$ holds for {every} onset (not one), and these mechanisms place every run in a $\psi$-location throughout $\kI$ up to any $\varphi$-release, $\rho$ holds for every $\nu\in Z_s$, matching Definition~\ref{def:sat}; no boundary is violated and (unlike a shared-clock treatment) no later onset's window is dropped.

Since $\varphi\in\Gamma_{s_0}$, the claim gives $\mathcal{T},(s_0,\nu_0)\models\varphi$. Therefore $\varphi$ is satisfiable.
\end{proof}




\subsection{Proof of Theorem~\ref{thm:completeness} -- Completeness of the tableau procedure}


\begin{proof}
Assume $\mathcal{T},(s_0,\nu_0)\models\varphi$ with $\mathcal{T}$ non-Zeno. We build a marking $\mathcal{M}$ and attach to each node $N=\langle\Gamma,\mathcal{X},Z,\obligations,\eventualities\rangle$ a configuration $(s,\nu)$ of $\mathcal{T}$ with $\nu\in Z$, maintaining the invariant: $(s,\nu)\models\gamma$ for all $\gamma\in\Gamma$, and every discrete successor of $(s,\nu)$ used in the construction satisfies $\mathcal{X}$. The root maps to $(s_0,\nu_0)$, where the invariant holds by assumption. It is preserved rule by rule.

{Boolean and context.} Rules~\eqref{tag:and},~\eqref{tag:ax} preserve the invariant by the semantics of $\land$ and of $\AX$ (the latter moving $\psi$ into $\mathcal{X}$, to be discharged at the successors). For~\eqref{tag:or} and the until/release splits, $(s,\nu)$ satisfies the principal formula, hence at least one disjunct; we select that OR-child. For until we take the stop child iff $\mathcal{T}$ fulfils the eventuality at $(s,\nu)$ -- in the timed case at an elapsed time in $\kI$, keeping $Z\cap(x_e\in\kI)$ non-empty -- and otherwise the go child. For release we select the partition region ($Z^{<\kI}$, $Z^{\in\kI}$, $Z^{>\kI}$, or early discharge) containing $\nu(x_\rho)$ and matching $\mathcal{T}$'s safety/release status; since the partition is disjoint and exhaustive, exactly one region applies and the chosen intersection is non-empty.

{Successor generation (AND-node).} For {each} $\EX\psi\in\Gamma$, the model has a discrete successor $(s',\nu')$ with $(s',\nu')\models\psi$, and, as $\mathcal{T}$ satisfies the $\AX$-obligations, $(s',\nu')\models\mathcal{X}$. We attach {all} these successor children to $\mathcal{M}$ and continue the construction at each from its model successor; the invariant is preserved at every child, and universal obligations in $\mathcal{X}$ hold at each because they hold at every successor of $(s,\nu)$. This is where the construction genuinely branches, in step with $\mathcal{T}$.

{Timed untils (dedicated checks).} A timed until $\xi=\tEU{\varphi}{\psi}{\kI}$ demanded at a node is verified by its check~\eqref{tag:chk}, which succeeds iff $Z\subseteq\mathrm{Sat}(\xi)$; this is where full completeness is secured, including under persistent contexts that re-demand $\xi$ at a continuum of onsets. We show $Z\subseteq\mathrm{Sat}(\xi)$. The node's zone $Z$ contains only valuations $\nu$ realised by the guiding model -- each is some reachable configuration $(s',\nu')$ of $\mathcal{T}$ at which $\xi$ is demanded, hence $\mathcal{T},(s',\nu')\models\xi$. So $\mathcal{T}$ supplies, {from that very valuation}, a witness path reaching $\psi$ within $\kI$ ($\varphi$ before, every safety obligation of $\mathcal{X}$ respected, since $\mathcal{T}\models\varphi$). The check, computing $\mathrm{Sat}(\xi)$ over $\mathcal{M}$ along $\mathcal{T}$'s own (possibly branching) witnesses, therefore places {every} $\nu\in Z$ in $\mathrm{Sat}(\xi)$ -- distinct onsets get their own witnesses, with no shared-window over-constraint, the defect of an accumulating scheme. Hence $Z\subseteq\mathrm{Sat}(\xi)$ and the check succeeds; the universal $\mathbf{A}$-variant is identical with $\mathcal{T}$'s witnesses ranging over all runs, which it does since $\mathcal{T}\models\tAU{\varphi}{\psi}{\kI}$ there. A demanded timed {release} is identical via its check~\eqref{tag:chkr}: $\mathcal{T}$ satisfies $\rho$ at every demanding configuration, so every $\nu\in Z$ lies in the greatest-fixpoint $\mathrm{Sat}(\rho)$ and $Z\subseteq\mathrm{Sat}(\rho)$ holds. Finiteness of $\mathrm{Sat}(\cdot)$ as a finite union of zones (the dense onsets are symbolic, never enumerated) follows from $\mathcal{T}$'s finite region quotient.

No selected node is closed: $\Gamma$ is propositionally consistent (it holds at $(s,\nu)$) and every chosen zone is non-empty (it contains $\nu$). Since only finitely many distinct nodes exist (Lemma~\ref{lem:finiteness}), every infinite path of $\mathcal{M}$ repeats a node, giving a candidate loop on that path; we check acceptance of each. For~(i): if some eventuality $e$ were postponed forever on a loop, the invariant would keep a model configuration satisfying $e$ at each node, but in TCTL a satisfied until is fulfilled after finite elapsed time (Definition~\ref{def:sat}); the model therefore forces a stop branch within the loop, removing $e$ -- contradiction. So every eventuality is fulfilled. For~(ii): the corresponding run of $\mathcal{T}$ is non-Zeno, so by the necessity direction of Lemma~\ref{lem:nonzeno} the loop is time-progressing (a never-reset clock kept below $M$ would cap the accumulated time and make the run Zeno). Hence every loop is accepting, $\mathcal{M}$ is a successful marking, and the tableau succeeds on $\varphi$.
\end{proof}

